\scanpage{57}
\mspage{13}

On va peut-être se borner pour la suite à des opérations de gommage de repères, ou de surmarquage, de nature spéciale, qui se traduisent, sur \inserted{chacun} des torseurs $R_a$ $(a\in A)$ concernés sous $\mathbb Z_a/\unsure{r(a)|\#(a)}$ \deleted{$\unsure{\mathbb Z/r(a)}$} $(\mathbb Z_a=\mathbb Z\mathbin{\underset{\{\pm1\}}{\wedge}}\omega(a))$, par une opération de \unsure{restriction|contraction} du groupe d'opérateurs de $\mathbb Z_a/r(a)$ à $\mathbb Z_a/r'(a)$, où $r'(a)\mid r(a)$ (opération de \uline{gommage} \uline{partiel} \inserted{\uline{uniforme}} en $a$) ou d'extension du groupe d'opérateurs de $\mathbb Z_a/r(a)$ à $\mathbb Z_a/r''(a)$, où $r''(a)$ \inserted{$=\unsure{d_{(a)}r(a)|d(a)r(a)}$} est un multiple de $r(a)$ (opération de \uline{surmarquage uniforme} en $a$, \deleted{\unsure{)}} de multiplicateur $d(a)$). On se rend compte \inserted{ici} que le gommage \uuline{partiel} est en fait un renforcement de structure, pas un affaiblissement, et inversement pour \unsure{le surmarquage uniforme|les surmarquages uniformes} --- qui, au niveau combinatoire, sera défini sans donnée de structure supplémentaire (sauf bien sûr les multiplicateurs $d(a)$).
% archive sheet number: 56

